% Generated draft: requires mathematical and source review.
\documentclass[11pt]{article}
\usepackage[margin=25mm]{geometry}
\usepackage{fontspec}
\setmainfont{Latin Modern Roman}
\usepackage{amsmath,amssymb,mathtools,unicode-math}
\setmathfont{Latin Modern Math}
\usepackage{enumitem,needspace,xurl,hyperref}
\setlength{\parindent}{0pt}
\setlength{\parskip}{5pt}
\setlength{\emergencystretch}{3em}
\newcommand{\N}{\mathbb{N}}
\newcommand{\Z}{\mathbb{Z}}
\newcommand{\Q}{\mathbb{Q}}
\newcommand{\R}{\mathbb{R}}
\newcommand{\C}{\mathbb{C}}
\newcommand{\F}{\mathbb{F}}
\newcommand{\A}{\mathbb{A}}
\newcommand{\PP}{\mathbb{P}}
\newcommand{\E}{\mathbb{E}}
\newcommand{\eps}{\varepsilon}
\newcommand{\dd}{\,\mathrm{d}}
\begin{document}
\section*{UnsolvedMath ID: 1359}
\textbf{Source problem code:} catalogue entry
\subsection{Definitions and mathematical statement}
Use the standard $9\times9$ Sudoku setting specified by the supplied review. Let $D=\{1,\ldots,9\}^2$ be the set of cells and let $\Sigma=\{1,\ldots,9\}$ be the symbols. A completed grid is a map $g:D\to\Sigma$ whose restriction to each row, each column, and each block
\[\{3a+1,3a+2,3a+3\}\times\{3b+1,3b+2,3b+3\},\qquad a,b\in\{0,1,2\},\]
is a bijection onto $\Sigma$. A puzzle is a partial map $P:C\to\Sigma$ with $C\subseteq D$; its entries are the clues. Let $\mathcal S(P)$ be the set of completed grids extending $P$.

A puzzle is minimal with a unique solution if
\[|\mathcal S(P)|=1\quad\text{and}\quad \left|\mathcal S\!\left(P|_{C\setminus\{c\}}\right)\right|\geq2\quad\text{for every }c\in C.\]
Here minimal means that no individual clue can be removed while preserving uniqueness; it does not mean that the number of clues is globally smallest. Write $\mathcal M$ for the finite set of these puzzles. The problem asks for an exact enumeration, but the record does not say whether the answer is $|\mathcal M|$ for fixed labelled cells and symbols, or the number of equivalence classes under a specified collection of symbol relabellings and grid symmetries. No equivalence relation is fixed by the source, so those counts must be distinguished before assigning one numerical target.

The background reports examples with 17 clues. The supplied literature review checked 2026-08-17 reports that no exhaustive exact count was verified, and quotes a statistical estimate of approximately $3.10\times10^{37}$ minimal puzzles. That estimate is not an exact count and its equivalence convention is not specified in the supplied record. The registry labels the problem open; the dated triage labels it partially solved / PARTIAL-PROGRESS.
\subsection{Short English statement}
Count standard $9\times9$ Sudoku puzzles that have exactly one completion and lose uniqueness whenever any clue is removed. An exact count requires fixing whether puzzles related by relabelling or grid symmetries are counted separately; the supplied numerical figure is only a statistical estimate.
\subsection{Sources}
\begin{itemize}
\item Local UnsolvedMath registry, permanent ID 1359, source code catalogue entry, ``Sudoku: Minimal Puzzles Count''; supplied statement and background.
\item ``Sudoku'' overview, accessed 2026-08-17 according to the supplied secondary reference: \url{https://en.wikipedia.org/wiki/Sudoku}. The registry review cites it for the absence of a precise count and the statistical estimate; no original enumeration or estimation paper is supplied.
\item Supplied literature triage checked 2026-08-17, PARTIAL-PROGRESS; it requests an equivalence convention and exact enumeration, and cautions that dated triage does not exclude later or unindexed work.
\end{itemize}
\end{document}
